% !TeX root = ../papers/nc-p3.tex
\section{非超楕円な種数三の滑らかな六次点グラフの排除}
\label{sec:nonhyperelliptic-sextic-exclusion}
\begin{theorem}\label{thm:nonhyperelliptic-sextic-exclusion}
有限Frobenius候補で一辺の点グラフが滑らかな非超楕円曲線 $C$（種数3、次数6）そのもので、
四辺の各射影ファイバーが空か一点となるものは存在しない。
\end{theorem}
\begin{proof}[点曲線の三次式]
各位置で次数6の偏極 $L$ は $h^0=4$ なので完全線型系となる。
$T=L/\omega_C$ は次数2で非有効である。有効ならその長さ2の因子が $L$ の切断に
一条件しか課さず、閉埋込みに反する。$\omega_C/T$ もRiemann--Rochで非有効である。
この空間曲線は二次曲面に入らない。滑らかな二次曲面なら二次数は $(2,4)$ または逆順となり
次数2の鉛筆ができる。二次錐の解消での類は $as+6f$、随伴種数は $-a^2+6a-5$。
種数3なら $a=2,4$ のうち負切断との交叉 $6-2a\ge0$ により $a=2$ となり同様である。

反対位置を変数にした前節の行列式三次式 $S_p$ が各 $p$ で非零であることを示す。
全関係空間が可分解元を含まないのは、非退化な整曲線上の非零切断の積が非零であるためである。
従って三次縮約の中支持は3以上（前述の特異な二元行列束の議論）である。
中支持4では定理\ref{thm:finite-cubic-separators}のUFD上の証明をそのまま使える。
中支持3で行列式が恒等零なら、定理\ref{thm:ternary-point-subalgebras}による
三文字の点グラフが平面であり、ある辺で共通三文字の可換関係を含む。
その三文字の零点を両側で $a,b\in C$ とする。両側の三座標の比が共通だから、
基点を除いた偏極 $L_i(-a)=L_{i+1}(-b)=N$ は同型である。
$a=b$ なら元の二偏極も一致し、曲線を含む二次曲面がないため六関係は交換子全体となって矛盾する。
$a\ne b$ なら次数7の $N(a+b)$ は非常に豊富で、五次元の全切断を持つ。
共通三文字とそれぞれの補文字 $x,y$ をこの標準でない射影四空間の座標とする。
残る三関係はそこで三つの二次式となる。共通三文字が非零の開部分ではこの三式の解が
そのまま二辺座標の対に復元されるため、零点は与えた曲線に限る。
外の割線に $xy$ の係数を持つ式がなければ $(j_i(a),j_{i+1}(b))$ が点対応に入ってしまう。
従って割線にも二指定点しかない。三式は完全交叉で次数8なのに、
台が次数7の整曲線だけとなり矛盾する。

反対位置の $L=L_2$ に戻る。$I_C$ は四つの三次式、三つの一次係数syzygyを持つ。
実際、二次制限は次元10の同型である。$L$ の基点自由鉛筆の完全列と
$n\ge3$ での $H^1(L^{n-2})=0$ から全高次の制限も全射になる。
従って曲線は算術的Cohen--Macaulayであり、$I_C$ は3正則である
（$H^1(I_C(2))=H^2(I_C(1))=H^3(I_C)=0$）。
Hilbert--BurchとHilbert関数により、四つの生成元は線型な $3\times4$ 行列 $A(z)$ の最大小行列である。
$A$ の曲線上の階数は常に2である。階数1以下では全小行列の微分が零となって滑らかさに反する。
\end{proof}

\begin{proof}[各長さ二の階数低下因子が三割線に乗ること]
定理\ref{thm:opposite-curve-divisor}の $D_p$ は長さ2であり、そのスキーム全体が $S_p$ の特異点に入る。
実際、三次行列式の微分は二次小行列の一次結合で、それら全てが $D_p$ 上で零である。
四つの符号付き最大小行列の列を $m$ として、$S_p=\ell m$、$\ell\ne0$ と書く。
曲線上で $dm$ の像は $\ker A$ に等しい。$Am=0$ を微分すれば包含が従い、
両階数が2であることは滑らかな曲線のイデアルを四式が生成することによる。
従って長さ2の $D_p$ 上では階数2の束 $\ker A$ が定数三平面 $W=\ker\ell$ に入る。
$W\to H^0(D_p,(W\otimes\mathcal O)/\ker A)$ の像は高々二次元である。
非零な定数核 $y$ を取れば、一次式の列 $A(z)y$ は $D_p$ とその線型包の直線で零になる。
この直線では $A$ の四つの最大小行列が一本の三次式の倍になり、その式は恒等零ではない
（$C$ は直線を含まない）。よって直線と $C$ の交叉は $D_p$ を含む長さ3の因子である。

$p$ を動かすとこの直線は正則に動く。長さ2の因子への切断評価が非常に豊富性で全射だからである。
曲線との交叉は有限で全ファイバー長が3、滑らかな基曲線上で平坦となる。
曲面 $C\times C$ 上のこの族は純な因子で、そこから $D$ を除いた残りは次数1のグラフ $r(p)$ である。
長さ3の交叉因子 $D_p+r(p)$ に対して $L(-D_p-r(p))$ は直線を含む平面から二切断を持つ。
Riemann--Rochにより
\[
 D_p+r(p)\ \sim\ T+s(p)
\]
となる一意な点 $s(p)$ がある。これは次数1のAbel埋込みを逆に読んだ射である。
どの点 $z$ でも $T+z$（加法的束記法）はただ一つの有効因子を持つ。
実際、$h^0(\omega_C/T)=0$ である。従って $s$ が定数なら $D$ の第二射影は有限台となり不可能である。
$r$ が定数 $r_0$ なら、全 $s(p)$ について $T+s(p)-r_0$ が有効である。
次数2のSerre双対により $s(p)$ は $(\omega_C/T)+r_0$ の唯一の有効因子に入ってしまう。
よって $r,s$ はともに非定数で、Hurwitzにより自己同型である。

もし $r(p)=s(p)$ なら $D_p\sim T$ となって非有効性に反する。
従って $g=rs^{-1}$ は固定点を持たない。$\tau=\sigma^{-1}$ として前節の因子の類を使い、
$p$ を動かすとJacobianに作用する準同型に次の恒等式を得る。
\[
 (1-g_*)s_* =1+\tau_* .
\]
\end{proof}

\begin{proof}[Fermat四次曲線への限定]
有限個のデータを複素数へ埋めて議論できる。標準平面上で固定点を持たない $g$ は、
相異なる三固有方向を持つ。重複なら固定直線が曲線と交わるからである。
固有座標 $u,v,w$ の純四乗の係数はいずれも非零であり、従って三固有値の四乗が等しい。
共通倍率と順序により $g$ は $(1,\mathrm i,-\mathrm i)$ に規格化できる。
これはJacobianの接空間上の実際の作用でもある。一般に平面四次式 $f$ に
$f(Mz)=c f(z)$ となる行列 $M$ が作用すると、標準微分の留数表示から接空間上の行列は
$(\det M/c)M$ である。ここで $c=\det M=1$ である。
同じ事実により、実際の接空間行列 $N$ については $f(Nz)=\det(N) f(z)$ となる。

有限自己同型群の作用を不変Hermite内積でユニタリとし、上のJacobian等式に
その随伴を右から掛けると $\tau+\tau^{-1}=-(g+g^{-1})$ である（接空間上の記号）。
従って $ku$ では $\tau=-1$、残る二平面では固有値が $\pm\mathrm i$ である。
両方の符号が出ると $\det\tau=-1$ となり、$u^4$ の非零係数に反する。
ゆえに残る二平面では一つのスカラー $\eta=\pm\mathrm i$ である。
$g$ 不変の四次単項式は $u^4,v^4,w^4,u^2vw,v^2w^2$ であり、
$\tau$ の行列式は1なので $u^2vw$ は現れない。
さらに $(1-g)s=1+\tau$ から、$s$ の $v,w$ の行は対角で、
対角成分の対は $(\mathrm i,1)$ または $(1,-\mathrm i)$ である。
ユニタリ性で $u$ 方向も保つ。純四乗の係数は保たれるが $v^2w^2$ の符号は変わるのでこれも消える。
従って $C$ はFermat平面四次曲線で、$\sigma$ は一座標を原始四乗根倍する射影変換である。
\end{proof}

\begin{proof}[偏極の全列挙と三次重なりによる排除]
$K=L_{i+2}L_i/L_{i+1}^2$ とする。$\sigma$ の四固定点は $u=0$ の標準超平面切断 $B_0,\ldots,B_3$ である。
$K\ne1$ なら $D$ は対角を含まない。含めば残りが負の自己交叉のグラフと数値同値なので
そのグラフに等しく、$K=1$ となる。$D\cdot\Delta=0$ ゆえ両者は交わらない。
対角への制限と $\mathcal O(\Delta)|_\Delta=\omega_C^{-1}$、
$\mathcal O(\Gamma_{\sigma^{-1}})|_\Delta=\mathcal O(\sum B_j)=\omega_C$ から $K^2=1$ を得る。
$K=1$ でも同じである。

$\sigma$ に不変な $\operatorname{Pic}^0$ の群は $(\mathbb Z/4)^2$ である。
実際、接空間の固有値から $1-\sigma^*$ の核の位数は
$|2(1-\eta)^2|^2=16$。各 $4B_j$ は $\omega_C$ と同値であり、
$A=B_1-B_0,B=B_2-B_0$ はこの核で独立な四捩れである。独立性が破れれば
$2A,2B,2(A-B)$ のいずれかが零となって次数2の鉛筆ができる。
$\sum B_j\sim4B_0$ より $B_3-B_0=-A-B$ である。
一階差 $d_i=L_{i+1}/L_i$ は $\sigma^*d_i=d_i K^4=d_i$ を満たすので $d_i^4=1$、
従って $\sigma^*L_i=L_i d_i^4 K^6=L_i$ となる。

不変な次数2の束を $2B_0+mA+nB$ と書く。不変かつ有効なら、
有効因子が唯一であり、その台は四固定点に限られる（それ以外では軌道の長さが4）。
従って有効な十個は $B_j+B_k$、残る非有効なものの $(m,n)$ は
\[
 (1,2),(1,3),(2,1),(2,3),(3,1),(3,2).
\]
各 $L_i/\omega_C$ はこのリストの点であり、位置 $i$ でのラベルは
$e+i d+\binom{i}{2}k$、$2k=0$ となる。隣接ラベルが等しいと二次曲面の非存在で可換となる。
これを除くと順序付き四列は30個である。

\path{computations/verify_nonhyperelliptic_sextic.py}で全30列の三次重なりを厳密に除外する。
モデル $u^4=t^4-1$ とし、$B_j$ の基点を $1,\mathrm i,-1,-\mathrm i$ の順に取る。
$L=\sum a_j B_j$ に $a=(6-m-n,m,n,0)$ を使う。全切断の基底は
\[
 \frac{u^e t^h}{\prod_j(t-t_j)^{b_j}},\qquad
 b_j=\lfloor(a_j+e)/4\rfloor,\quad 0\le e<4,\quad 0\le h\le\sum b_j-e .
\]
分岐点での位数と無限遠での極の上限、巡回群の固有分解から従う。
各隣接対の16積を共通分母にすると、標数101、$\mathrm i=10$ で階数は10である。
従って標数零でも全六関係が積の核であり、良い素点での特殊化は計算した六平面になる。
三次交叉の次元は特殊化で下がらない。各始点で次元0の四列が6個、次元1が16個、次元2が8個であり、
いずれも標準有限Frobeniusに必要な次元4に達しない。以上で全候補を除外した。
\end{proof}
